\writeSectionFrame{\presentationTitle}{Memento}
\begin{frame}{Steckbrief}
	\begin{itemize}
 		\item Deutscher Name: Memento
 		\item Alternativer Name: Token
 		\item Englischer Name: Memento
 		\item Gruppe: Verhaltensmuster
 	\end{itemize}
\end{frame}

\begin{frame}{Memento}
	\begin{block}{Zweck}
 		Erfasse und externalisiere den internen Zustand eines Objekts, ohne seine Kapselung zu verletzen, so dass das Objekt in diesen Zustand zurückgesetzt werden kann. 
 	\end{block}
\end{frame}

\begin{frame}{Allgemeines Klassendiagramm}
    \begin{tikzpicture}[scale=0.7, transform shape]
        % Define the UML classes
        \umlclass[x=0, y=0]{Originator}{
            - state: State\\
        }{
            + createMemento(): Memento\\
            + setMemento(m: Memento)
        }
    
        \umlclass[x=6, y=0]{Memento}{
            - state: State\\
        }{
            + getState(): State
        }
    
        \umlclass[x=3, y=-3]{Caretaker}{
            - memento: Memento\\
        }{
            + saveMemento()\\
            + undo()
        }
    
        \umluniassoc{Originator}{Memento}
        \umlaggreg{Caretaker}{Memento}
    \end{tikzpicture}
\end{frame}
\note{
	Der Urheber ist eine beliebige Klasse, deren Zustand ganz oder zum Teil gespeichert werden soll. Das ist die Aufgabe der Klasse Memento. Der Urheber kann ein Memento erzeugen und zurückgeben. Der Aufbewahrer bewahrt ein Memento	oder mehrere auf. Allerdings soll er den aufbewahrten Zustand nicht einsehen oder verändern können. Mit der Methode setzeMemento() des Urhebers kann ein	gespeicherter Schnappschuss zurückgesetzt werden. Der Aufbewahrer fordert den	Urheber später auf, in einen früheren Zustand zurückzukehren.
}

\frameWithImageOnly{\BILDERPFAD/memento1.jpg}
\note{
	Das Beispiel simuliert ein Zeichnungsprogramm, bei dem der Zustand einer Zeichnung (mit mehreren Formen) gespeichert und wiederhergestellt werden kann. 
}

\begin{frame}{Klassendiagramm - Zeichenprogramm}
    \begin{tikzpicture}[scale=0.7, transform shape]
        \umlclass[x=0, y=0]{Originator}{
            -shapes:List<String>
        }{
            +removeShape(shape:String)\\
            +printShapes() \\
            +saveToMemento(): Memento \\
            +restoreFromMemento(memento:Memento)
        }
    
        \umlclass[x=8, y=0]{Memento}{
            - shapes: List<String>\\
        }{
        }

        \umlclass[x=3, y=-5]{Caretaker}{
            - mementoStack: Stack<Memento>\\
        }{
            + undo(originator: Originator)\\
            + saveState(originator: Originator)\\
        }
    
        \umluniassoc{Originator}{Memento}
        \umlaggreg{Caretaker}{Memento}
    \end{tikzpicture}
\end{frame}

\begin{frame}[fragile]{Urheber}
	\begin{exampleblock}{Originator.java}
 		\begin{lstlisting}
public class Originator {
    private final List<String> shapes;

    public void addShape(String shape) {
        shapes.add(shape);
    }

    public Memento saveToMemento() {
        return new Memento(shapes);
    }

    public void restoreFromMemento(Memento memento) {
        shapes.clear();
        shapes.addAll(memento.getShapes());
    }
}

 		\end{lstlisting}
 	\end{exampleblock}
\end{frame}

\begin{frame}[fragile]{Memento}
	\begin{exampleblock}{Memento.java}
 		\begin{lstlisting}
public class Memento {
    private final List<String> shapes;

    public Memento(List<String> shapes) {
        this.shapes = new ArrayList<>(shapes);
    }

    public List<String> getShapes() {
        return shapes;
    }
}
 		\end{lstlisting}
 	\end{exampleblock}
\end{frame}

\begin{frame}[fragile]{Aufbewahrer}
	\begin{exampleblock}{Caretaker.java}
 		\begin{lstlisting}
public class Caretaker {
    private final Stack<Memento> mementoStack = new Stack<>();

    public void saveState(Originator originator) {
        mementoStack.push(originator.saveToMemento());
    }

    public void undo(Originator originator) {
        if (!mementoStack.isEmpty()) {
            originator.restoreFromMemento(mementoStack.pop());
        } else {
            System.out.println("...");
        }
    }
}
 		\end{lstlisting}
 	\end{exampleblock}
\end{frame}

\frameWithImageOnly{\BILDERPFAD/memento2.jpg}
\note{
	Im Folgenden simulieren wir ein ganz einfaches Versionisierungssystem. Wir legen Textdateien an, speichern den Stand als Commit und können die Versionen wieder zurückspielen.  
}

\begin{frame}{Klassendiagramm - Mini-Git}
    \begin{tikzpicture}[scale=0.9, transform shape]
        % Define the UML classes
        \umlclass[x=0, y=0]{TextFile}{
        }{}
    
        \umlclass[x=6, y=0]{TextFileMemento}{
        }{}
    
        \umlclass[x=3, y=-3]{TextFileVersionManager}{
        }{}
    
        \umluniassoc{TextFile}{TextFileMemento}
        \umlaggreg{TextFileVersionManager}{TextFileMemento}
    \end{tikzpicture}
\end{frame}

\begin{frame}[fragile]{Urheber}
	\begin{exampleblock}{TextFile.java}
 		\begin{lstlisting}
public class TextFile {
	private int id;
	private String content;
	private String name;
	
    public TextFileMemento createSnapShot(int serialNo, 
            String commitMessage){
        return new TextFileMemento(serialNo, content,
            commitMessage, new Date());
    }

    public void writeContent(String updatedContent) {
        this.content = updatedContent;
    }
}
 		\end{lstlisting}
 	\end{exampleblock}
\end{frame}

\begin{frame}[fragile]{Memento}
	\begin{exampleblock}{TextFileMemento.java}
 		\begin{lstlisting}
public class TextFileMemento {
	private final int serialNumber;
	private final String content;
	private final String commitMessage;
	private final Date commitTime;
	
	public TextFileMemento(int serialNumber, String content, 
            String commitMessage, Date commitTime) {
		this.serialNumber = serialNumber;
		this.content = content;
		this.commitMessage = commitMessage;
		this.commitTime = commitTime;
	}

	public int getSerialNumber() {
		return serialNumber;
	}
}
 		\end{lstlisting}
 	\end{exampleblock}
\end{frame}

\begin{frame}[fragile]{Aufbewahrer}
	\begin{exampleblock}{TextFileVersionManager.java}
 		\begin{lstlisting}
public class TextFileVersionManager {
	private Map<Integer, List<TextFileMemento>> states;
	
    public void commit(TextFile file, String commitMessage) {
        List<TextFileMemento> snapshots = 
            states.get(file.getId());
        if (snapshots == null) {
            snapshots = new ArrayList<>();
            states.put(file.getId(), snapshots);
        }

        states.get(file.getId()).add(file.createSnapShot(
        		snapshots.size(), commitMessage));

        System.out.println("...");
    }
 		\end{lstlisting}
 	\end{exampleblock}
\end{frame}

\begin{frame}[fragile]{Aufbewahrer}
	\begin{exampleblock}{TextFileVersionManager.java}
 		\begin{lstlisting}
public void listCommits(TextFile file) {
    List<TextFileMemento> commits = states.get(file.getId());
    if (commits == null || commits.isEmpty()) {
        System.out.println("...");
    } else {

        System.out.println("...");
        for (TextFileMemento commit : commits) {
            System.out.println(commit);
        }
    }
}
 		\end{lstlisting}
 	\end{exampleblock}
\end{frame}

\begin{frame}[fragile]{Aufbewahrer}
	\begin{exampleblock}{TextFileVersionManager.java}
 		\begin{lstlisting}
public void restore(TextFile file, int commitSerialNo) {
    List<TextFileMemento> commits = states.get(file.getId());
    if (commits != null && commitSerialNo > -1) {
        TextFileMemento versionToRestore = 
            commits.get(commitSerialNo);
        file.writeContent(versionToRestore.getContent());

        states.put(file.getId(), commits.subList(0, 
            commitSerialNo));

        System.out.println("Rollback to: " + commitSerialNo);
    }
}
 		\end{lstlisting}
 	\end{exampleblock}
\end{frame}